\chapter{Metodología y planificación}
\label{chap:metodologia}

\lettrine{E}{l} desarrollo de PadelCenter se ha llevado a cabo de forma individual,
sin un equipo que requiera ceremonias de coordinación, pero sí con una metodología
de trabajo deliberada que ha guiado el orden de las tareas, la gestión del riesgo
y el seguimiento del progreso. En este capítulo se describe dicha metodología, la
planificación inicial del trabajo y un análisis de cómo la ejecución real se
desvió de lo previsto.

\section{Metodología de desarrollo}
\label{sec:metodologia_desarrollo}

Se adoptó un enfoque \textbf{iterativo e incremental por módulo funcional},
combinado con el carácter \textit{contract-first} ya descrito en el
capítulo~\ref{chap:diseno}. Cada módulo (autenticación, centros y pistas,
reservas, torneos) se desarrolló siguiendo el mismo ciclo:

\begin{enumerate}
  \item Definición del contrato del módulo en la especificación OpenAPI.
  \item Diseño del modelo de dominio y los puertos correspondientes.
  \item Implementación de los casos de uso y los adaptadores de infraestructura.
  \item Cobertura con tests unitarios y de integración.
  \item Implementación del frontend que consume el módulo.
\end{enumerate}

Este ciclo se eligió frente a un desarrollo en cascada porque permite obtener,
al final de cada iteración, un incremento funcional completo y verificable
(\textit{vertical slice}), reduciendo el riesgo de descubrir problemas de
integración entre capas al final del proyecto. Frente a un marco más formal
como Scrum, se descartaron las ceremonias de equipo (\textit{daily},
\textit{sprint planning} con varios roles, retrospectivas grupales) por no
aportar valor en un trabajo individual, conservando en cambio los elementos que
sí son útiles en solitario: tablero de tareas, iteraciones acotadas en el tiempo
y revisión del progreso al cierre de cada una.

\subsection{Herramientas de seguimiento}

\begin{itemize}
  \item \textbf{Control de versiones y tablero de tareas}: Git, con un
    repositorio independiente para backend y frontend, e issues para registrar
    tareas pendientes y errores detectados durante el desarrollo.
  \item \textbf{Integración continua}: ejecución automática de la batería de
    tests (unitarios, integración y arquitectura) en cada cambio, evitando que
    se acumule deuda técnica no detectada.
  \item \textbf{Especificación OpenAPI}: además de contrato de API, ha actuado
    como documento vivo de seguimiento del alcance funcional: el progreso de
    cada módulo es directamente visible en el número de endpoints especificados
    e implementados.
\end{itemize}

\section{Plan de trabajo inicial}
\label{sec:plan_inicial}

La planificación inicial dividió el trabajo en cinco fases, alineadas con la
estructura de la memoria:

\begin{table}[ht]
\centering
\begin{tabularx}{\textwidth}{lXc}
\hline
\textbf{Fase} & \textbf{Contenido} & \textbf{Esfuerzo estimado} \\
\hline
F1 — Análisis y diseño & Estado del arte, definición de requisitos, diseño de
  arquitectura y modelo de dominio, primera versión del contrato OpenAPI &
  20\% \\
F2 — Backend & Implementación de dominio, casos de uso, persistencia,
  seguridad y tests del backend & 40\% \\
F3 — Frontend & Implementación de la interfaz, integración con la API
  generada, gestión de estado y autenticación en cliente & 25\% \\
F4 — Pruebas e integración & Tests de integración end-to-end, pulido de
  errores de integración entre frontend y backend & 10\% \\
F5 — Documentación & Redacción de la memoria y preparación de la defensa &
  5\% \\
\hline
\end{tabularx}
\caption{Plan de trabajo inicial por fases, con esfuerzo estimado como
  porcentaje del total.}
\label{tab:plan_inicial}
\end{table}

\subsection{Análisis de riesgos}

\begin{longtable}{p{3.4cm}p{3.9cm}p{4.8cm}}
\caption{Riesgos identificados al inicio del proyecto y mitigaciones aplicadas.}
\label{tab:riesgos} \\
\hline
\textbf{Riesgo} & \textbf{Impacto} & \textbf{Mitigación adoptada} \\
\hline
\endfirsthead
\multicolumn{3}{c}{\tablename\ \thetable{} -- \textit{(continúa de la página anterior)}} \\
\hline
\textbf{Riesgo} & \textbf{Impacto} & \textbf{Mitigación adoptada} \\
\hline
\endhead
\hline
\multicolumn{3}{r}{\textit{Continúa en la página siguiente}} \\
\endfoot
\hline
\endlastfoot
Disponibilidad irregular por solapamiento con la actividad laboral del autor &
  Retrasos en la cadencia de avance, especialmente durante periodos de
  mayor carga laboral &
  Priorización estricta del alcance funcional (sección~\ref{sec:objetivos})
  frente a funcionalidades secundarias \\
Cambios en el contrato de la API tras iniciar la implementación del frontend &
  Retrabajo en backend y frontend simultáneamente &
  Enfoque \textit{contract-first}: cualquier cambio se refleja primero en el
  OpenAPI y se propaga mediante generación de código, evitando
  inconsistencias manuales \\
Subestimación de la complejidad del módulo de torneos (generación de
  emparejamientos) &
  Retraso en la fase de backend &
  Implementación incremental por formato de competición, empezando por el más
  sencillo (\texttt{ROUND\_ROBIN}) antes de abordar eliminación y grupos \\
Curva de aprendizaje de herramientas no usadas previamente por el autor
  (Orval, ArchUnit, Testcontainers) &
  Retraso puntual al inicio de cada módulo nuevo &
  Prueba de concepto aislada de cada herramienta antes de integrarla en el
  flujo principal de desarrollo \\
Condiciones de carrera al reservar la misma pista y franja horaria desde
  varios clientes simultáneamente &
  Doble reserva sobre el mismo recurso e inconsistencia de datos &
  Comprobación de solapamiento y creación de la reserva en la misma
  transacción de base de datos, respaldada por una restricción de unicidad
  que garantiza atomicidad aunque dos peticiones lleguen a la vez \\
Elección de una arquitectura monolítica en lugar de microservicios,
  con el consiguiente riesgo de comprometer la escalabilidad y el
  aislamiento de fallos a medida que creciera el sistema &
  Escalado todo-o-nada, despliegues acoplados entre módulos y ausencia de
  aislamiento ante un fallo grave de un módulo &
  Monolito modular sobre arquitectura hexagonal: las fronteras entre
  \textit{centros}, \textit{reservas}, \textit{torneos} y \textit{usuarios}
  quedan explícitas desde el dominio, de forma que una futura extracción a
  microservicios no exigiría reescribir la lógica de negocio \\
\end{longtable}

\section{Ejecución real y desviaciones}
\label{sec:ejecucion_real}

La figura~\ref{tab:ejecucion_real} resume la actividad real de desarrollo por
mes, medida a partir del historial de control de versiones de ambos
repositorios.

\begin{table}[ht]
\centering
\begin{tabularx}{\textwidth}{lXX}
\hline
\textbf{Periodo} & \textbf{Backend} & \textbf{Frontend} \\
\hline
Noviembre 2025 -- febrero 2026 &
  Actividad baja y discontinua: diseño de la arquitectura hexagonal, modelo de
  dominio inicial y primera versión del contrato OpenAPI, en paralelo con
  la actividad laboral del autor &
  Sin actividad: el frontend no se inició hasta tener estabilizado el
  contrato de la API \\
Marzo 2026 &
  Sin actividad registrada &
  Sin actividad registrada \\
Abril 2026 &
  Fase intensiva: implementación de la mayor parte de los casos de uso,
  persistencia, seguridad y tests &
  Fase intensiva: implementación de la mayor parte de las páginas, integración
  con la API generada y gestión de estado \\
Mayo -- junio 2026 &
  Cierre, correcciones puntuales y redacción de la memoria &
  Cierre, correcciones puntuales y redacción de la memoria \\
\hline
\end{tabularx}
\caption{Actividad real de desarrollo por periodo, según el historial de
  control de versiones.}
\label{tab:ejecucion_real}
\end{table}

La desviación más relevante respecto al plan inicial es la \textbf{concentración
del trabajo en un periodo mucho más corto del previsto}: la fase de análisis y
diseño (F1) se extendió de forma discontinua durante varios meses por
solapamiento con las obligaciones laborales del autor, mientras que las
fases F2 y F3 (backend y frontend), planificadas inicialmente de forma
secuencial, terminaron ejecutándose de manera prácticamente simultánea durante
abril de 2026 una vez estabilizado el contrato de la API. Esto fue posible
precisamente gracias al enfoque \textit{contract-first}: al existir un contrato
OpenAPI estable, ambos proyectos pudieron avanzar en paralelo sin bloquearse
mutuamente, ya que el frontend podía generar sus tipos y \textit{hooks} a partir
de la especificación sin depender de que el backend estuviese completamente
implementado.

En términos de esfuerzo relativo, la proporción final entre backend y frontend
(78 \textit{commits} frente a 20) es más desequilibrada que la estimada
inicialmente en la tabla~\ref{tab:plan_inicial} (40\,\% frente a 25\,\%). Esto
se explica por el mayor número de módulos de dominio a modelar en el backend
(siete entidades principales, 42 casos de uso) frente a un frontend que reutiliza
en gran medida los mismos patrones de página para cada módulo una vez establecida
la estructura inicial de rutas y componentes.
